\chapter{برنامه‌نویسی یعنی چه؟}%
\label{ch:what-is-programming}

پیش از آنکه حتی یک خط کد بنویسیم، باید بدانیم اصلاً «برنامه» چیست و رایانه
چگونه آن را می‌فهمد. این فصل هیچ کدی ندارد؛ فقط چند ایدهٔ ساده اما بنیادی
دارد که تمامِ کتاب روی آن‌ها بنا می‌شود.

\section{رایانه: ماشینی سریع اما ساده‌دل}

بسیاری گمان می‌کنند رایانه باهوش است. حقیقت برعکس است: رایانه به‌خودیِ‌خود
هیچ‌چیز نمی‌فهمد. نه منظورِ شما را حدس می‌زند، نه اشتباهتان را می‌بخشد و نه
ابتکاری از خود نشان می‌دهد. تنها هنرش این است که دستورهای بسیار ساده را با
سرعتی باورنکردنی (میلیاردها دستور در ثانیه) و بدونِ خستگی اجرا می‌کند.

پس اگر رایانه این‌قدر ساده‌دل است، این همه کارِ شگفت‌انگیز (بازی‌ها،
پیام‌رسان‌ها، نقشه‌ها) از کجا می‌آید؟ پاسخ: از \emph{برنامه‌نویس‌ها}؛
انسان‌هایی که کارهای بزرگ را به دستورهای کوچکِ قابلِ فهم برای ماشین
شکسته‌اند.

\begin{note}
این جمله را به خاطر بسپارید، چون بارها به دردتان می‌خورد: \emph{رایانه دقیقاً
همان کاری را می‌کند که به او گفته‌اید؛ نه آن کاری که می‌خواستید بکند.} اگر
برنامه‌ای درست کار نکرد، مشکل تقریباً همیشه در دستورهایی است که ما داده‌ایم.
\end{note}

\section{برنامه چیست؟}

\textbf{برنامه} فهرستی از دستورهاست که به ترتیب نوشته شده‌اند تا رایانه با
اجرای آن‌ها کاری را انجام دهد. همین. یک بازیِ رایانه‌ایِ بزرگ هم در نهایت
همین است: میلیون‌ها دستورِ کوچک که با نظمی حساب‌شده کنارِ هم چیده شده‌اند.

با یک مثالِ روزمره شروع کنیم. فرض کنید می‌خواهید به دوستی که تا حالا چای دم
نکرده، چای دم کردن یاد بدهید. ممکن است بنویسید:

\begin{enumerate}
  \item کتری را از آب پُر کن.
  \item کتری را روشن کن.
  \item صبر کن تا آب بجوشد.
  \item چای را در قوری بریز.
  \item آبِ جوش را روی چای بریز.
  \item ده دقیقه صبر کن.
\end{enumerate}

این یک «برنامه» برای انسان است. برنامهٔ رایانه‌ای هم دقیقاً همین ساختار را
دارد: دنباله‌ای از گام‌های روشن و بدونِ ابهام. تفاوت فقط در مخاطب است؛
دوستِ شما اگر بنویسید «چای دم کن» خودش بقیه را می‌فهمد، اما رایانه
\emph{هیچ‌چیز} را خودش نمی‌فهمد و باید همه‌چیز را مو‌به‌مو گفت.

\section{زبانِ برنامه‌نویسی چیست؟}

دستورها را به چه زبانی بنویسیم؟ زبانِ روزمرهٔ ما پر از ابهام است. اگر بگوییم
«یک لیوان آب بریز»، منظورمان دقیقاً چقدر آب است؟ تا لبِ لیوان؟ نصفه؟ رایانه
با ابهام میانه‌ای ندارد. به همین دلیل، زبان‌هایی مخصوصِ گفت‌وگو با رایانه
ساخته شده‌اند: \textbf{زبان‌های برنامه‌نویسی}. این زبان‌ها واژگانی اندک اما
قاعده‌هایی بسیار دقیق دارند، طوری که هر جمله فقط و فقط یک معنی داشته باشد.

در این کتاب از زبانِ \textbf{سلام} استفاده می‌کنیم که دستورهایش را می‌توان به
فارسی نوشت. برای نمونه، جملهٔ «عبارتِ سلام دنیا را نمایش بده» در زبانِ سلام
چنین نوشته می‌شود: \code{چاپ "سلام دنیا"}. کوتاه، دقیق و بدونِ ابهام.

\section{مترجمی به نامِ کامپایلر}

یک نکتهٔ فنیِ کوچک باقی می‌ماند. قلبِ رایانه (پردازنده) حتی زبانِ سلام را هم
مستقیم نمی‌فهمد؛ او فقط رشته‌هایی از صفر و یک را اجرا می‌کند که به آن
\emph{زبانِ ماشین} می‌گویند. پس به مترجمی نیاز داریم که نوشتهٔ ما را به زبانِ
ماشین برگرداند. این مترجم \textbf{کامپایلر} نام دارد. کامپایلرِ زبانِ سلام
برنامه‌ای است به نامِ \code{salam} که در فصلِ \ref{ch:first-program} با آن
کار خواهیم کرد.

کامپایلر یک لطفِ بزرگِ دیگر هم در حقِ ما می‌کند: متنِ برنامه را وارسی
می‌کند و اگر جایی از قاعده‌های زبان تخطی کرده باشیم، با پیامی روشن به ما
\textbf{خطا} می‌دهد. تازه‌کارها از دیدنِ خطا دلسرد می‌شوند؛ اشتباه است.
خطا دشمنِ شما نیست، معلمِ شماست: دارد می‌گوید دقیقاً کجا را باید درست کنید.

\begin{deep}[نگاهی به درون: چرا صفر و یک؟]
درونِ رایانه میلیاردها کلیدِ الکترونیکیِ بسیار ریز هست که هر کدام یا روشن‌اند
یا خاموش. روشن را «۱» و خاموش را «۰» می‌نامیم. هر داده‌ای (عدد، متن، عکس،
صدا) در نهایت با همین صفر و یک‌ها نمایش داده می‌شود و هر دستوری هم الگویی از
همین صفر و یک‌هاست. خوشبختانه شما هرگز مجبور نیستید صفر و یک بنویسید؛
کامپایلر این ترجمه را برایتان انجام می‌دهد.
\end{deep}

\section{برنامه‌نویس چه‌کار می‌کند؟}

شاید فکر کنید کارِ اصلیِ برنامه‌نویس «تایپ کردنِ کد» است. در واقع بخشِ
کوچکی از کارِ اوست. کارِ اصلیِ برنامه‌نویس \emph{فکر کردن} است:

\begin{enumerate}
  \item \textbf{فهمیدنِ مسئله:} دقیقاً چه می‌خواهیم؟ ورودی چیست و خروجی چه
        باید باشد؟
  \item \textbf{طراحیِ راهِ‌حل:} مسئله را به گام‌های کوچک و روشن می‌شکنیم.
        به این گام‌ها \emph{الگوریتم} می‌گویند؛ موضوعِ فصلِ بعد.
  \item \textbf{نوشتنِ برنامه:} الگوریتم را با زبانِ برنامه‌نویسی
        می‌نویسیم.
  \item \textbf{آزمودن و اصلاح:} برنامه را اجرا می‌کنیم، خطاها را می‌یابیم و
        درستشان می‌کنیم. این چرخه بارها تکرار می‌شود و کاملاً طبیعی است.
\end{enumerate}

این کتاب هر چهار مهارت را کنارِ هم پرورش می‌دهد، اما بیشترین تأکیدش بر
گام‌های ۱ و ۲ است؛ چون تایپ کردن را همه بلدند، \emph{فکر کردنِ ساخت‌یافته}
است که برنامه‌نویس می‌سازد.

\section{تمرین‌ها}

\exercise سه دستگاهِ اطرافِ خود را نام ببرید که درونشان برنامهٔ رایانه‌ای
اجرا می‌شود (راهنمایی: فقط به «رایانه» فکر نکنید؛ ماشینِ لباس‌شویی چطور؟).

\exercise «دندان مسواک زدن» را مانندِ مثالِ چای، در پنج تا هشت گامِ دقیق
بنویسید. سپس آن را به دوستی بدهید و بپرسید آیا جایی از آن مبهم است.

\exercise فرض کنید به رایانه گفته‌اید «۱۰ بار کلمهٔ سلام را چاپ کن» ولی
منظورتان ۵ بار بوده است. رایانه چند بار چاپ می‌کند؟ چرا این پرسش مهم است؟
